% Generated draft: requires mathematical and source review.
\documentclass[11pt]{article}
\usepackage[margin=25mm]{geometry}
\usepackage{fontspec}
\setmainfont{Latin Modern Roman}
\usepackage{amsmath,amssymb,mathtools,unicode-math}
\setmathfont{Latin Modern Math}
\usepackage{enumitem,needspace,xurl,hyperref}
\setlength{\parindent}{0pt}
\setlength{\parskip}{5pt}
\setlength{\emergencystretch}{3em}
\newcommand{\N}{\mathbb{N}}
\newcommand{\Z}{\mathbb{Z}}
\newcommand{\Q}{\mathbb{Q}}
\newcommand{\R}{\mathbb{R}}
\newcommand{\C}{\mathbb{C}}
\newcommand{\F}{\mathbb{F}}
\newcommand{\A}{\mathbb{A}}
\newcommand{\PP}{\mathbb{P}}
\newcommand{\E}{\mathbb{E}}
\newcommand{\eps}{\varepsilon}
\newcommand{\dd}{\,\mathrm{d}}
\begin{document}
\section*{UnsolvedMath ID: 800128}
\textbf{Source problem code:} catalogue entry
\subsection{Definitions and mathematical statement}
Under the conventional interpretation, let \(G=\mathrm{PL}_c(\R^2)\) be the group, under composition, of piecewise-linear homeomorphisms \(f:\R^2\to\R^2\) with compact support, where
\[
\operatorname{supp}(f)=\overline{\{x\in\R^2:f(x)\ne x\}}.
\]
Regard \(G\) as a discrete group, and use trivial actions on the coefficient groups \(\Z\) and \(\R\). The supplied excerpt does not explicitly define its convention for \(\mathrm{PL}_c\); this conventional interpretation requires confirmation against the source.

A real group \(3\)-cocycle is a function \(c:G^3\to\R\) satisfying, for all \(g_1,g_2,g_3,g_4\in G\),
\[
\begin{aligned}
0={}&c(g_2,g_3,g_4)-c(g_1g_2,g_3,g_4)
+c(g_1,g_2g_3,g_4)\\
&-c(g_1,g_2,g_3g_4)+c(g_1,g_2,g_3).
\end{aligned}
\]
Its class \([c]\in H^3(G;\R)\) defines an evaluation homomorphism
\[
\Phi_c:H_3(G;\Z)\longrightarrow(\R,+),\qquad
\alpha\longmapsto\langle[c],\alpha\rangle.
\]
The explicit question is whether there exists such a cocycle \(c\) for which \(\Phi_c\) is surjective; equivalently, whether
\[
\exists[c]\in H^3(G;\R)\quad\forall r\in\R\quad
\exists\alpha\in H_3(G;\Z):\ \langle[c],\alpha\rangle=r.
\]
The source frames this as a particular instance of asking whether an analog of discrete \(\mathrm{GV}\) exists in higher dimensions. Neither discrete \(\mathrm{GV}\) nor the requirements for such an analog are defined in the supplied record, so that broader question remains unspecified beyond the displayed cohomological target.
\subsection{Short English statement}
Does the discrete group of compactly supported piecewise-linear homeomorphisms of \(\R^2\) admit a real group \(3\)-cocycle whose evaluation on integral third homology attains every real number? The source presents this as a higher-dimensional analog of discrete \(\mathrm{GV}\), without defining that broader requirement.
\subsection{Sources}
\begin{itemize}
\item Local UnsolvedMath registry: permanent ID \(800128\), exact source problem code \texttt{catalogue entry}, titled ``Foliations and Diffeomorphism Groups --- Question 69.'' The supplied record attributes the problem to Nariman and Yazdi, with proposed year \(2025\), and labels it open. This status is supplied registry metadata, not an independently verified mathematical assessment.
\item Nariman, Yazdi, \emph{Problem List on Foliations and Diffeomorphism Groups} (2025), Question 69. Supplied source URL: \url{https://arxiv.org/abs/2503.18274}. The record reports access on \texttt{2026-10-16} and extraction by \texttt{arxiv\_latex}; these are supplied provenance details, not a claim of independent source access or verification.
\end{itemize}
\end{document}
